% Generated draft: requires mathematical and source review.
\documentclass[11pt]{article}
\usepackage[margin=25mm]{geometry}
\usepackage{fontspec}
\setmainfont{Latin Modern Roman}
\usepackage{amsmath,amssymb,mathtools,unicode-math}
\setmathfont{Latin Modern Math}
\usepackage{enumitem,needspace,xurl,hyperref}
\setlength{\parindent}{0pt}
\setlength{\parskip}{5pt}
\setlength{\emergencystretch}{3em}
\newcommand{\N}{\mathbb{N}}
\newcommand{\Z}{\mathbb{Z}}
\newcommand{\Q}{\mathbb{Q}}
\newcommand{\R}{\mathbb{R}}
\newcommand{\C}{\mathbb{C}}
\newcommand{\F}{\mathbb{F}}
\newcommand{\A}{\mathbb{A}}
\newcommand{\PP}{\mathbb{P}}
\newcommand{\E}{\mathbb{E}}
\newcommand{\eps}{\varepsilon}
\newcommand{\dd}{\,\mathrm{d}}
\begin{document}
\section*{UnsolvedMath ID: 10400056}
\textbf{Source problem code:} catalogue entry
\subsection{Definitions and mathematical statement}
The Kontsevich integral is a diagram-valued universal finite-type invariant, customarily written for a knot as
\[
 Z(K)\in\widehat{\mathcal A}(S^1)=\prod_{d\geq0}\mathcal A^{(d)}(S^1).
\]
The completion records all homogeneous degrees, and the diagram relations and normalization must match the intended invariant.

J. Roberts asks for a \emph{good topological construction} of this integral. A construction of the intended invariant should specify its topological input and diagram-valued output and justify agreement with the chosen Kontsevich integral, including any framing or normalization corrections. Merely restating that the existing invariant is invariant under isotopy does not define what new construction or explanation the question seeks.

The source does not give a formal meaning of ``good'' or ``topological'', nor specify whether the desired input category consists of knots, links, framed tangles or another setting. It does not explicitly forbid every analytic ingredient, require a purely combinatorial algorithm, or prescribe tangle functoriality. Those stronger requirements are not added. The summary notes configuration-space integral constructions as existing context but supplies no precise criterion distinguishing an acceptable new topological construction.
\subsection{Short English statement}
Give a satisfactory topological construction of the Kontsevich integral, as requested by J. Roberts. The source leaves the intended topological mechanism, input category, normalization, and criterion of goodness unspecified.
\subsection{Sources}
\begin{itemize}
\item Local UnsolvedMath registry, permanent ID 10400056, source code catalogue entry; statement attribution J. Roberts; proposer field and proposal year not supplied; update 2026-08-06.
\item Ohtsuki, \emph{Problems on invariants of knots and 3-manifolds} (2002), Problem 3.6, PDF page 54: \url{https://msp.org/gtm/2002/04/gtm-2002-04-024s.pdf}. Registry access date 2026-07-29; PDF font-structure extraction. Provenance distinguishes 231 numbered bold headings. Present primary-literature status and redistribution terms require review.
\item The supplied summary mentions existing configuration-space integral constructions but gives no exact additional reference or verified answer to Roberts's formulation.
\end{itemize}
\end{document}
